\documentclass{sbs-exam}

\graphicspath{
  {../Assets/}
  {../Media/}
}

% Override \MarksBreakdown to accept our 8-argument version
\makeatletter
\let\orig@MarksBreakdown\MarksBreakdown
\def\MarksBreakdown#1#2#3#4#5#6#7#8{%
  \gdef\@MCQ{#1}\gdef\@MCQNo{#2}\gdef\@TFQ{#3}\gdef\@TFQNo{#4}\gdef\@SAQ{#5}\gdef\@SAQNo{#6}\gdef\@LAQTotal{#7}\gdef\@TotalMarks{#8}%
}
\makeatother

% Override marksheet to match our actual structure
\makeatletter
\def\makemarksheet{%
  \begin{center}
  \begin{tabular}{|p{0.5\textwidth}|p{0.35\textwidth}|}
    \hline
    \multicolumn{1}{|l|}{\bfseries Mark Breakdown:} & \multirow{3}{*}{Student Name:} \\
    \cline{1-1}
    Part I: Multiple Choice Questions $\@MCQ \times \@MCQNo$ & \\
    Part II: True/False Questions $\@TFQ \times \@TFQNo$ & \\
    Part III: Fill in the Blanks $\@SAQ \times \@SAQNo$ & \\
    Part IV: Long Answer Questions $\@LAQTotal$ & \\
    \cline{1-1}
    \multicolumn{1}{|l|}{Total Marks: \@TotalMarks} & Score: \fillline \\
    \hline
  \end{tabular}
  \end{center}
}
\makeatother

\examtitle{G10 A3 Mathematics Re-Midterm, 2025-2026 Term 2}
\examsubject{G10 A3 Mathematics Pure 1 \& 2}{Apr 2026}{60 Minutes}{Shi Feng}
\MarksBreakdown{2 marks}{10}{2 marks}{10}{2 marks}{10}{30 marks}{90 marks}

% Section heading helpers (used in template)
\makeatletter
\def\@LAQTotal{30 marks}
\def\imarks#1{\hfill\textbf{(#1 marks)}}
\makeatother

\begin{document}

\makemarksheet

\examnotice

\examtools

\newpage

\makeatletter
\section*{Multiple Choice Questions (\@MCQ $\times$ \@MCQNo)}
\makeatother
\begin{enumerate}
    \item Simplify $\dfrac{x^2 + 7x + 12}{x+3}$.
    \begin{enumerate}
        \item $x+4$
        \item $x-4$
        \item $x+3$
        \item $x-3$
    \end{enumerate}

    \item Find the remainder when $x^3 - 3x^2 + 2x - 1$ is divided by $(x-2)$.
    \begin{enumerate}
        \item $-1$
        \item $1$
        \item $3$
        \item $-3$
    \end{enumerate}

    \item Which of the following is a factor of $x^3 - 4x^2 - 3x + 18$?
    \begin{enumerate}
        \item $x+2$
        \item $x-2$
        \item $x+3$
        \item $x-3$
    \end{enumerate}

    \item The midpoint of the line segment joining $(-2,5)$ and $(4,-1)$ is
    \begin{enumerate}
        \item $(1,2)$
        \item $(2,1)$
        \item $(1,1)$
        \item $(2,2)$
    \end{enumerate}

    \item The equation of the circle with centre $(-1,2)$ and radius $3$ is
    \begin{enumerate}
        \item $(x+1)^2 + (y-2)^2 = 9$
        \item $(x-1)^2 + (y+2)^2 = 9$
        \item $(x+1)^2 + (y-2)^2 = 3$
        \item $(x-1)^2 + (y+2)^2 = 3$
    \end{enumerate}

    \item $\log_2 64$ is equal to
    \begin{enumerate}
        \item $5$
        \item $6$
        \item $8$
        \item $32$
    \end{enumerate}

    \item Solve $3^x = 81$.
    \begin{enumerate}
        \item $2$
        \item $3$
        \item $4$
        \item $5$
    \end{enumerate}

    \item The discriminant of the quadratic $x^2 - 3x + 4 = 0$ is
    \begin{enumerate}
        \item $-7$
        \item $7$
        \item $-5$
        \item $5$
    \end{enumerate}

    \item The solutions to $x^2 - 7x + 10 = 0$ are
    \begin{enumerate}
        \item $2$ and $5$
        \item $-2$ and $-5$
        \item $1$ and $10$
        \item $-1$ and $-10$
    \end{enumerate}

    \item The gradient of the line perpendicular to $y = -2x + 5$ is
    \begin{enumerate}
        \item $2$
        \item $-2$
        \item $\dfrac{1}{2}$
        \item $-\dfrac{1}{2}$
    \end{enumerate}
\end{enumerate}

\makeatletter
\section*{True or False Questions (\@SAQ $\times$ \@SAQNo)}
\makeatother
State whether each statement is True or False.

\begin{enumerate}
    \item The factor theorem states that if $f(p)=0$ then $(x-p)$ is a factor of $f(x)$. \hfill True \quad False

    \item For the circle $(x-h)^2 + (y-k)^2 = r^2$, the centre is $(h,k)$ and radius is $r$. \hfill True \quad False

    \item $\log_a x - \log_a y = \log_a (x-y)$. \hfill True \quad False

    \item The quadratic equation $x^2 + 2x + 3 = 0$ has two distinct real roots. \hfill True \quad False

    \item The remainder theorem says that when $f(x)$ is divided by $(ax-b)$, the remainder is $f\!\left(\dfrac{b}{a}\right)$. \hfill True \quad False

    \item The graph of $y = 2^x$ is an increasing function. \hfill True \quad False

    \item The perpendicular bisector of a chord of a circle always passes through the centre. \hfill True \quad False

    \item $\log_a 1 = 0$ for any $a>0$, $a\neq 1$. \hfill True \quad False

    \item Completing the square for $x^2 - 8x$ gives $(x-4)^2 - 16$. \hfill True \quad False

    \item The line $y = x + 3$ is a tangent to the circle $(x-1)^2 + (y+1)^2 = 8$. \hfill True \quad False
\end{enumerate}

\makeatletter
\section*{Fill in the Blanks (\@SAQ $\times$ \@SAQNo)}
\makeatother
Complete each statement with the correct answer.

\begin{enumerate}
    \item When simplifying the algebraic fraction $\dfrac{x^2+5x+4}{x+1}$, after cancelling common factors, the result is \fillline.

    \item The remainder when $x^3-2x^2+3x-4$ is divided by $(x-1)$ is \fillline.

    \item If $(x+2)$ is a factor of $x^3+ax^2-3x-2$, then the value of $a$ is \fillline.

    \item The midpoint of the line segment joining $A(2,-3)$ and $B(-4,5)$ is \fillline.

    \item The equation of the circle with centre $(-3,1)$ and radius $2$ is $(x+3)^2+(y-1)^2 = \fillline$.

    \item The value of $\log_4 64$ is \fillline.

    \item Solve $2^{3x-1}=32$; then $x = \fillline$.

    \item The discriminant of the quadratic equation $3x^2-2x+1=0$ is \fillline.

    \item Completing the square for $x^2+10x$ gives $(x+5)^2 - \fillline$.

    \item A circle has equation $x^2+y^2+4x-6y-3=0$. The $y$-coordinate of its centre is \fillline.
\end{enumerate}

\newpage

\section*{Long Answer Questions}

\begin{enumerate}

\item % Q1: Polynomials and factor theorem
The polynomial $f(x) = 2x^3 - 3x^2 - 8x + 3$.
\begin{enumerate}
    \item Use the factor theorem to show that $(x+1)$ is a factor of $f(x)$. \imarks{3}
    \item Hence factorise $f(x)$ completely into linear factors. \imarks{4}
    \item Write down all the real roots of the equation $f(x)=0$. \imarks{3}
\end{enumerate}
\vspace{6cm}

\item % Q2: Quadratics
Consider the quadratic function $f(x) = x^2 - 4x - 5$.
\begin{enumerate}
    \item Solve the equation $f(x)=0$. \imarks{2}
    \item Write $f(x)$ in the form $(x-p)^2 + q$. \imarks{3}
    \item Hence state the minimum value of $f(x)$ and the value of $x$ at which it occurs. \imarks{2}
    \item Find the discriminant of $f(x)$ and explain what it tells you about the roots. \imarks{3}
\end{enumerate}
\vspace{8cm}

\item % Q3: Coordinate geometry
The circle $C$ has equation $(x+2)^2 + (y-3)^2 = 25$.
\begin{enumerate}
    \item Write down the centre and radius of $C$. \imarks{2}
    \item Verify that the point $P(1,-1)$ lies on $C$. \imarks{2}
    \item Find the equation of the tangent to $C$ at $P$. Give your answer in the form $ax+by+c=0$, where $a$, $b$ and $c$ are integers. \imarks{4}
\end{enumerate}
\vspace{6cm}

\item % Q4: Exponentials and logarithms
\begin{enumerate}
    \item Solve $4^{2x-1} = 64$. \imarks{3}
    \item Solve $\log_3(x+2) + \log_3(x-2) = 2$. \imarks{4}
    \item Given that $\log_a 25 = 2$, find the value of $a$. \imarks{3}
\end{enumerate}
\vspace{6cm}

\end{enumerate}

\end{document}
